problem:  
Let \((M^{2n+1}, \theta)\) be a compact, strictly pseudoconvex, simply connected, pseudo-Einstein pseudohermitian manifold of dimension \(2n+1 \ge 5\). Assume:  
1. the Tanaka–Webster scalar curvature \(R_\theta > 0\) everywhere, and  
2. the CR Paneitz operator \(P_\theta\) is **nonnegative** in the quadratic form sense,  
\[
\int_M (P_\theta f)\, f\, dV_\theta \ge 0 \quad \text{for all smooth } f.
\]  

Consider the CR \(Q'\)-curvature associated to \(P_\theta\), the pseudohermitian analogue of Branson’s \(Q\)-curvature. Suppose further that \(\int_M Q'_\theta\, dV_\theta > 0\).  

**Question:**  
Does it follow that \((M,\theta)\) is globally CR equivalent (up to scaling of \(\theta\)) to the standard pseudohermitian sphere \(S^{2n+1}\subset \mathbb{C}^{n+1}\)?  

Equivalently, do there exist simply connected, closed, pseudo-Einstein CR manifolds with \(R_\theta > 0\), nonnegative \(P_\theta\), and positive total \(Q'_\theta\) that are not globally spherical?  

---

Why is it a "good" problem:  

1. **Conceptual depth and novelty.**  
   This problem asks whether combined positivity of both second-order and fourth-order CR curvature invariants (\(R_\theta\) and \(P_\theta\), \(Q'_\theta\)) enforces geometric rigidity. While the CR Obata–Jerison–Lee theorem classifies constant Webster scalar curvature solutions in the CR Yamabe setting, no rigidity criterion involving the fourth-order Paneitz operator or its curvature analogue \(Q'_\theta\) is known. The open status of global results under \(P_\theta \ge 0\) and \(Q'_\theta > 0\) makes this a genuinely unresolved and forward-looking question.

2. **Bridge between analytic and geometric aspects.**  
   The problem sits at the confluence of **subelliptic PDE**, **global CR geometry**, and **analysis on the Heisenberg group**. It seeks CR analogues of conformal rigidity theorems that connect positivity of Paneitz/Q-curvature to sphericity or conformal flatness. Solving it may require extending analytic techniques (spectral properties of \(P_\theta\), compactness under CR Yamabe flow, or CR GJMS theory) into a genuinely pseudohermitian context. This integration of analytic, geometric, and topological tools is characteristic of impactful modern research problems.

3. **Connection with known results yet a genuine extension.**  
   Existing work by Graham–Lee, Case–Chang–Yang, and Hirachi establishes local positivity of the CR Paneitz operator only under perturbative assumptions or near-sphericity, and deals mostly with stability results—**not with global rigidity** under positivity of \(P_\theta\) and \(Q'_\theta\). By extending to a global geometric question and assuming simple connectivity but not explicit near-sphericity, this problem goes beyond the current frontier.

4. **Extreme simplicity with profound implications.**  
   The question can be stated concisely as:  
   > “Do positivity of \(R_\theta\), \(P_\theta\), and \(Q'_\theta\) force the CR sphere?”  
   Its formulation is transparent, yet it encodes deep subtleties involving the hierarchy of CR invariants. A positive solution would create a fourth-order CR analogue to Obata-type rigidity, while a counterexample would reveal new global phenomena in pseudohermitian geometry.  

Thus, this problem stands as a clean, plausible, and truly open research-level question that could shape the frontier of CR geometric analysis.